\documentclass{report}
\usepackage{amsmath,amssymb}
\setlength{\parindent}{0mm}
\setlength{\parskip}{1em}
\begin{document}
\begin{center}
\rule{6in}{1pt} \
{\large Boyana Norris \\
{\bf Automatic Kernel Acceleration of PETSc Krylov Solvers}}

MCS Division \\ Bldg 240 \\ Rm 2152 \\ Argonne National Laboratory \\ 9700 South Cass Avenue \\ Argonne \\ IL 60439
\\
{\tt norris@mcs.anl.gov}\\
Cherkuri Choudary\\
Deepan Balasubramanian\\
Jeswin Godwin\\
Daniel Lowell\\
Azamat Mametjanov\\
Boyana Norris\\
	P. Sadayappan, Gerald Sabin, Sravya Tirukkovalur\end{center}

In the last few years, the use of accelerator architecture in both
desktop and high-performance platforms has been rapidly increasing.
Cray’s XK6 systems have a hybrid architecture using NVIDIA Graphics
Processing Units (GPUs), and in a similar spirit, Intel has created the
Many Integrated Core (MIC) architecture codenamed Knights Corner.
Developing software for such heterogeneous architectures has become
increasingly important, but requires significant expertise and
development effort. Furthermore, the code developed and optimized for one
accelerator or MIC architecture is generally not easily portable to
another. One approach that addresses this challenge is to define code
transformations that enable the same reference implementation of an
algorithm to be ported to multiple platforms. An extension of this
approach is to enable users to express an algorithm in a high-level
language, which can then be transformed automatically into multiple
low-level source code representations for execution on different
accelerator architectures. Examples of low-level target codes include C
with OpenMP directives or CUDA C kernels.

Krylov subspace solvers are widely used iterative methods for solving
large, sparse, linear systems of equations, with implementations provided
by a number of popular numerical toolkits such as the Portable Extensible
Toolkit for Scientific Computation (PETSc). PETSc’s open-source,
extensible, and modular framework allows computational scientists or
other library developers to create new or modify existing solvers as
needed.

The implementation of a Krylov solver for execution on a GPU platform
takes place within the framework of PETSc vector and matrix data types.
PETSc has several GPU implementations of iterative solvers such as the
Generalized Minimum Residual Method (GMRES). These implementations use
PETSc’s vector and matrix data structures, but implement the algorithms
by calling GPU libraries external to PETSc. These libraries are provided
by different CUDA toolkits (e.g., the CUSP and THRUST libraries). The
delegation of control to GPU libraries provides a modular design which
hides the lower level implementations on the GPU, leaving the algorithm
developer free to use the familiar high-level operations. This black-box
approach is often suboptimal because the presence of multiple library
calls hampers locality-exploiting optimizations of low-level kernels. Our
transparent implementation of vector and matrix types keep intact the
modular design of PETSc while allowing the developer full access to the
tuning of low-level operations.

We explored a number of optimizations in our vector and matrix data
structures, including GPU global memory alignment, reduction of data
transfers between the host and the device, and maximizing GPU
utilization. Examples of these optimizations include asynchronous
streaming kernels and memory calls, pipelined reductions of large array
operations, use of dynamically allocated shared memory, register
offloading of global memory, and pinned host memory. The goal of this
exploratory step was to create the data structures and support functions
that can serve as templates for automated code generation and tuning of
the core kernel operations of a Krylov iterative solver.

After obtaining kernel templates, we began developing code
transformations of C-based implementations to CUDA C-based optimized
versions. These transformations were developed within the optimizing
autotuner Orio. Orio is a recently developed, extensible, and portable
software generation framework for empirical performance tuning. It takes
annotated C or Fortran source code as input, performs different
performance-tuning transformations on the internal code representation,
generates multiple code variants corresponding to the different
transformations, empirically evaluates the performance of the generated
codes, and selects the best-performing variant using several popular
heuristic search algorithms. Orio annotations consist of semantic
comments that specify the computation and optionally, the transformations
to perform. A separate tuning specification contains various parametrized
performance-tuning directives and bounds of optimization space to search.
In addition to the general-purpose tuning directives, such as loop fusion
and unrolling, tiling, and scalar replacement, Orio supports a number of
architecture-specific optimizations (e.g., generating calls to SIMD
intrinsics on Intel and Blue Gene/P architectures). Orio is implemented
in Python and can be dynamically extended by placing new language modules
(containing implementations of parsing, transformation, and code
generation interfaces) into the canonical path. Typically, most of the
execution time is spent in loops; hence, we have extended Orio’s Loop
module, which provides a C-like syntax for specifying array-based
computations. The extension parses user-provided GPU transformation
directives and generates CUDA kernels for execution on GPU devices. The
annotated loops are replaced by CUDA-specific resource marshaling for
parallel execution (e.g., data transfer, memory allocation, and thread
layout) and unmarshaling upon the conclusion of the GPU computation.

The proposed generation of kernels allows us to create a testbed for
generating and tuning iterative solver implementations and testing them
within the PETSc framework, including benchmarking the performance of new
GPU implementations against existing ones. We will present a detailed
comparison of the performance of our approach and existing manually-tuned
and library-based PETSc CPU and GPU implementations.

Future work includes enabling the developer to use vectors and (sparse)
matrices as base types and support code generation and tuning for a core
set of linear algebra operations. For example, instead of explicitly
iterating over index vector elements in, say an AXPY operation, one can
specify just y = y + alpha * x. Other future work includes implementing
more complex kernel code generation, integrating more aggressive
optimizations into the generated kernel variants, and accelerating other
solvers.


\end{document}
